\documentclass[conference]{IEEEtran}
\usepackage{xeCJK}
\setmainfont{Times New Roman}
\setsansfont{Arial}
\setmonofont{SimSun}
\setCJKmainfont{SimSun}
\setCJKsansfont{SimHei}
\setCJKmonofont{SimSun}
\xeCJKDeclareCharClass{CJK}{"2500->"257F,"2580->"259F}
\usepackage{booktabs}
\usepackage{array}
\usepackage{fancyvrb}
\usepackage{graphicx}
\usepackage{amsmath}
\usepackage{needspace}
\usepackage[hidelinks]{hyperref}
\renewcommand{\thesection}{\arabic{section}}
\renewcommand{\thesubsection}{\thesection.\arabic{subsection}}
\renewcommand{\thesubsubsection}{\thesubsection.\arabic{subsubsection}}
\renewcommand{\thesectiondis}{\arabic{section}}
\renewcommand{\thesubsectiondis}{\thesectiondis.\arabic{subsection}}
\renewcommand{\thesubsubsectiondis}{\thesubsectiondis.\arabic{subsubsection}}
\renewcommand{\thetable}{\arabic{table}}
\renewcommand{\thefigure}{\arabic{figure}}
\renewcommand{\tablename}{表}
\renewcommand{\figurename}{图}
\renewcommand{\IEEEkeywordsname}{关键词}
\renewcommand{\abstractname}{摘要}
\renewcommand{\appendixname}{附录}
\renewcommand{\arraystretch}{1.05}
\setlength{\tabcolsep}{4pt}
\makeatletter
\def\subsubsection{\@startsection{subsubsection}{3}{\z@}%
  {1.2ex plus .2ex minus .2ex}{0.5ex}%
  {\normalfont\normalsize\bfseries\raggedright}}
\makeatother
\begin{document}
\title{GPT-5.6-sol 在 WikiSQL 上的 Text-to-SQL 基线}
\author{\IEEEauthorblockN{Anonymous Author(s)}\IEEEauthorblockA{Affiliation withheld for anonymous review}}
\maketitle
\begin{center}\small
\textbf{实验日期：}2026 年 11 月 18 日\\
\textbf{数据集：}WikiSQL\\
\textbf{主模型：}gpt-5.6-sol（Responses API）\\
\textbf{报告类型：}Test 集正式实验结果 + 人工语义评测 pilot\\
\textbf{报告版本：}v3（汇入 20 条人工语义标注）
\end{center}
\vspace{2mm}
\begin{abstract}
本文系统评估大语言模型 gpt-5.6-sol 在 WikiSQL 单表 Text-to-SQL 任务上的零样本与少样本基线。主系统 E4 采用 Schema-aware 提示词、基于 Train 集的动态 5-shot 检索、无损表格序列化和 high reasoning effort，在完整 Test 集 15,878 条样本上取得 \textbf{85.63\% Logical Form Accuracy（LFA）} 与 \textbf{90.77\% Execution Accuracy（EX）}，95\% Wilson 置信区间分别为 [85.07\%, 86.16\%] 与 [90.31\%, 91.21\%]。全部预测通过 JSON 与结构化 SQL 校验，执行成功率为 99.94\%。
在同一前 1,000 条 Test 子集上的成对消融显示：动态 5-shot 检索相对 0-shot 提升 3.00 个百分点 LFA 与 3.40 个百分点 EX（McNemar 精确检验 p \textless{} 0.001）；单独开启 high reasoning effort 未带来可观测收益（E1 vs E2：p = 0.8145；E3 vs E4：p = 0.8145）。
在 Test 前 20 条 pilot 上进行的单标注人工语义评测中，20/20 被判为语义正确（95\% CI [83.89\%, 100\%]）；同一子集的自动 LFA/EX 均为 90\%，两条自动判错样本分别源于金标错误与问题歧义。该 pilot 仅用于暴露评价口径问题，不作为总体 Semantic Accuracy。
分层分析表明，系统在聚合查询、非等值条件、数值条件和宽表上明显退化：SUM 的 LFA 仅 44.12\%，COUNT 为 59.18\%；含 \texttt{\textgreater{}}/\texttt{\textless{}} 条件的样本 LFA 为 72.10\%，而全部为等值条件时为 88.69\%；含数值条件的样本 LFA 为 78.76\%，字符串条件为 86.52\%。此外，本文对数据规模、表长尾分布、提示词冻结、采样参数与成本口径进行了审计，发现现有实验在 Test 集复用、Dev 验证规模、提示词冻结证据链和缓存成本归因方面存在不足。本文认为该结果可作为可复核的工程基线，但尚不足以支撑“已完成系统性 Prompt 调优”或“SOTA”类结论。
\end{abstract}
\begin{IEEEkeywords}
Text-to-SQL；WikiSQL；大语言模型；少样本检索；执行准确率；实验审计
\end{IEEEkeywords}
\section{引言}
WikiSQL 是单表自然语言到 SQL 的代表性基准，其目标是在给定一张表和一句自然语言问题的条件下，生成 \texttt{SELECT} 列、聚合函数与 \texttt{WHERE} 条件组成的逻辑形式。与 Spider 等多表跨库任务相比，WikiSQL 的 SQL 空间受限，但由于真实表格存在列名歧义、取值多样、表规模长尾等问题，仍然是检验 schema linking、条件值匹配和聚合识别的重要基线。
本文关注的问题是：在不微调云端模型的前提下，仅依靠提示词、动态示例检索与结构化输出约束，当前大语言模型能在 WikiSQL 上达到什么水平；这些增益来自哪个组件；以及现有实验记录是否足以支撑相应结论。
本文的贡献如下：
\begin{enumerate}
\item 在完整 Test 集（15,878 条）上给出 gpt-5.6-sol 的可复核基线结果，并报告置信区间。
\item 在同一 1,000 条子集上完成 2×2 成对消融（0/5-shot × thinking off/on），并给出精确 McNemar 检验。
\item 补充按聚合类型、条件数量、检索分数、表行数、列数、算子类型和值类型的多维分层分析。
\item 从数据集、实验设计、算法与可复现性四个层面审计现有实验，明确指出结论的适用边界。
\end{enumerate}
\section{WikiSQL 数据集介绍与统计分析}
\subsection{数据集概览}
WikiSQL 由 Wikipedia 表格及人工标注的问题-SQL 对组成（Zhong et al., 2017；官方数据、评测脚本与说明见 \url{https://github.com/salesforce/WikiSQL} 。每条样本包含：
\begin{itemize}
\item 一张表：列名（header）、列类型（text / real）、若干行单元格；
\item 一个自然语言问题；
\item 一个逻辑形式：选择列 \texttt{sel}、聚合算子 \texttt{agg}、条件列表 \texttt{conds}，每个条件为 \texttt{[列索引, 算子, 值]}。
\end{itemize}
本地数据快照包含 Train / Dev / Test 三个划分，示例数分别为 56,355 / 8,421 / 15,878，合计 80,654 条。表文件共 26,531 张表（Train 18,585、Dev 2,716、Test 5,230），其中被问题引用的表分别为 17,984 / 2,630 / 5,069；其余为未被引用的表。三个划分的表 ID \textbf{完全互不重叠}（交集均为 0），这避免了跨划分的表级泄漏。
\begin{quote}\small
说明：公开文献对 WikiSQL 表数量的统计口径存在差异。本报告以本地数据文件为准。
\end{quote}
\begin{table*}[t]
\centering
\footnotesize
\caption{数据划分与表级统计}
\label{tab:t1}
\begin{tabular}{@{}lrrrrrrrrrr@{}}
\toprule
划分 & 示例数 & 表文件数 & 被引用表数 & 平均行数/表 & 行数中位数 & 行数 P95 & 最大行数 & \textgreater{}100 行的表 & 平均列数 & 最大列数 \\
\midrule
Train & 56,355 & 18,585 & 17,984 & 17.1 & 11 & 44 & 1,956 & 195 & 6.35 & 44 \\
Dev & 8,421 & 2,716 & 2,630 & 16.5 & 11 & 43 & 433 & 27 & 6.31 & 20 \\
Test & 15,878 & 5,230 & 5,069 & 18.4 & 12 & 48 & 1,950 & 75 & 6.32 & 23 \\
\bottomrule
\end{tabular}
\end{table*}
Test 集问题长度平均 11.7 个词，P95 为 20 个词，最大 43 个词；每个表平均被 3.04 个问题引用，最多 6 个。Test 集包含 15,832 个不同问题字符串，即存在 46 个重复问题。
\subsection{Test 集 SQL 结构分布}
\begin{table}[t]
\centering
\footnotesize
\caption{Test 集 SQL 结构分布（N = 15,878）}
\label{tab:t2}
\begin{tabular}{@{}llrr@{}}
\toprule
维度 & 类别 & 数量 & 占比 \\
\midrule
聚合 & NONE & 11,324 & 71.32\% \\
聚合 & MAX & 973 & 6.13\% \\
聚合 & MIN & 933 & 5.88\% \\
聚合 & COUNT & 1,455 & 9.16\% \\
聚合 & SUM & 587 & 3.70\% \\
聚合 & AVG & 606 & 3.82\% \\
条件数 & 0 & 131 & 0.83\% \\
条件数 & 1 & 10,832 & 68.22\% \\
条件数 & 2 & 3,879 & 24.43\% \\
条件数 & 3+ & 1,036 & 6.52\% \\
条件算子 & \texttt{=} & 18,114 & 82.96\% \\
条件算子 & \texttt{\textgreater{}} & 1,927 & 8.82\% \\
条件算子 & \texttt{\textless{}} & 1,805 & 8.27\% \\
条件值 & 字符串 & 16,424 & 75.19\% \\
条件值 & 数值 & 5,422 & 24.81\% \\
\bottomrule
\end{tabular}
\end{table}
可以看到三个显著的不均衡：
\begin{enumerate}
\item \textbf{聚合不均衡}：71.32\% 的查询不带聚合，而 COUNT/SUM/AVG 合计仅 16.68\%。微平均 LFA 会被 NONE 类主导。
\item \textbf{条件数集中}：68.22\% 的查询只有 1 个条件，2 个及以上条件占 30.96\%。
\item \textbf{算子与值类型偏斜}：等值条件占 82.96\%，字符串值占 75.19\%；\texttt{\textgreater{}}/\texttt{\textless{}} 与数值条件属于少数但更难的分支。
\end{enumerate}
\subsection{表规模长尾与上下文成本}
Test 集表行数呈明显长尾：75 张表超过 100 行（1.43\%），3 张超过 500 行，3 张超过 1,000 行，最大 1,950 行。由于 E4 采用“保留全部行”的无损表格序列化，长表会直接转化为输入 token：
\begin{itemize}
\item E4 平均输入 6,694 tokens/样本，P95 为 8,928，P99 为 12,739，最大 \textbf{132,634} tokens；
\item 输入超过 10k tokens 的样本 492 条（3.10\%），超过 20k 的 30 条（0.19\%），超过 30k 的 19 条（0.12\%）。
\end{itemize}
这意味着长表虽然比例低，但决定了上下文窗口、延迟和成本的上界，是该方法最脆弱的运行边界之一。
\subsection{数据集层面的任务特征}
\begin{itemize}
\item \textbf{单表、有限 SQL 语法}：不支持 JOIN、子查询、GROUP BY、HAVING、ORDER BY、多列聚合等，因此结果不能直接外推到 Spider 等复杂场景。
\item \textbf{模式简单但取值开放}：列名往往来自 Wikipedia 表头，存在缩写、同义与大小写差异；条件值可能带单位、百分号、日期或数字。
\item \textbf{存在逻辑形式等价}：不同 SQL 可能在当前表上得到相同结果，导致 EX 高于 LFA（本文 E4 相差 5.15 个百分点）。
\item \textbf{金标不含 \texttt{OP=3}}：本地 Train/Dev/Test 金标条件算子只有 \texttt{=}、\texttt{\textgreater{}}、\texttt{\textless{}}；但提示词允许模型输出算子 3，造成结构合法但执行器不支持的情况（E4 共 10 条）。
\end{itemize}
\section{基线算法与实验策略}
\subsection{总体框架}
\begin{figure*}[t]
\centering
\begin{Verbatim}[fontsize=\scriptsize]
问题 q + 目标表 T
        │
        ├── 动态 5-shot 检索（仅从 Train 表/问题/SQL 中检索）
        │        └── 词法相似度 + SQL 结构启发式 + 表头重合度
        │
        ├── Prompt 构造
        │        ├── Schema-aware developer 指令
        │        ├── few-shot 示例（问题 + 紧凑表 + 金标 SQL）
        │        └── 目标问题 + 无损紧凑表
        │
        ├── Responses API（strict JSON Schema）
        │        └── {sel, agg, conds}
        │
        └── 逻辑形式校验 → SQLite 执行 → LFA / EX
\end{Verbatim}
\end{figure*}
实现位于 \texttt{experiment/\allowbreak{}run\_modern\_baseline.py}
\subsection{模型与推理参数}
\needspace{3.6cm}
\noindent\begin{minipage}{\linewidth}
\centering
\scriptsize
\begin{tabular}{@{}>{\raggedright\arraybackslash}p{2.84cm}>{\raggedright\arraybackslash}p{4.56cm}@{}}
\toprule
项目 & 取值 \\
\midrule
模型 & \texttt{gpt-5.6-sol} \\
接口 & Responses API，\texttt{/\allowbreak{}v1/\allowbreak{}responses} \\
输出约束 & \texttt{text.format = json\_schema}，\texttt{strict = true} \\
输出上限（代码设定） & 256 tokens（E1/E3）、384 tokens（E2/E4）；不完整时回退 768 \\
reasoning effort & E1/E3 为 \texttt{none}，E2/E4 为 \texttt{high} \\
采样参数 & 未在代码中设置；原始响应显示服务端默认 \texttt{temperature=1.0}、\texttt{top\_p=0.98} \\
微调 & 无 \\
\bottomrule
\end{tabular}
\end{minipage}\par\smallskip
需要指出：实际记录中输出 tokens 最大值达到 1,117（E4），说明网关并未严格执行代码设定的 \texttt{max\_output\_tokens}；该约束与“不完整回退”逻辑在本轮实验中未真正发挥作用。
\subsection{Prompt 设计}
developer prompt 采用固定 Schema-aware 指令，要求模型：
\begin{enumerate}
\item 只输出 \texttt{sel}、\texttt{agg}、\texttt{conds}；
\item 列索引从 0 开始；
\item 聚合算子取 0–5；
\item 条件为 \texttt{[列索引, 整数算子, 值]}，不允许输出符号算子；
\item 只依据问题与表格作答；
\item 保留条件中的原始取值；
\item 不输出解释。
\end{enumerate}
user prompt 由“问题 + 紧凑表格 JSON”组成，表格包含列名、列类型与全部行。输出由 strict JSON Schema 约束，避免自由文本解析。
\begin{quote}\small
注：\textbf{Dev 验证不构成“调优”证据}：四个 Prompt 版本各只跑 2 条/组，且全部满分。更准确的表述是“Prompt 预先设计，Dev 仅做流程冒烟验证”。
\end{quote}
\subsection{动态少样本检索}
检索器仅使用 Train 集，评分由三部分构成：
\begin{itemize}
\item 问题与示例问题的词法相似度（IDF 加权、问题结构启发式）；
\item 目标表表头与示例表表头的重合度；
\item SQL 结构启发式（聚合、条件数量等）。
\end{itemize}
候选集先按词法重叠剪枝，再排序取 Top-5。每个示例的表被投射到其 SQL 使用到的列，但保留\textbf{全部行}；示例 SQL 使用 Train 金标。Test 金标不进入提示词。
\subsection{实验条件}
\begin{table*}[t]
\centering
\footnotesize
\caption{四组实验设计}
\label{tab:t4}
\begin{tabular}{@{}lrllrl@{}}
\toprule
实验 & Few-shot & Thinking & 推理强度 & Test 样本数 & 角色 \\
\midrule
E1 & 0-shot & Off & none & 1,000 & 基础对照 \\
E2 & 0-shot & On & high & 1,000 & Thinking 消融 \\
E3 & 5-shot & Off & none & 1,000 & Few-shot 消融 \\
E4 & 5-shot & On & high & 15,878 & 主基线 \\
\bottomrule
\end{tabular}
\end{table*}
E1–E3 为节省成本只使用 Test 前 1,000 条（非随机抽样）；E4 覆盖完整 Test。为公平比较，本文同时报告 E4 在同一前 1,000 条上的结果。\textbf{消融样本量小}：1,000 条子集足以做配对检验，但不能代表全 Test 分布
\subsection{评测指标}
\needspace{3.2cm}
\noindent\begin{minipage}{\linewidth}
\centering
\scriptsize
\begin{tabular}{@{}>{\raggedright\arraybackslash}p{3.24cm}>{\raggedright\arraybackslash}p{4.16cm}@{}}
\toprule
指标 & 定义 \\
\midrule
LFA（Logical Form Accuracy） & \texttt{sel} 与 \texttt{agg} 完全一致；\texttt{conds} 按无序集合比较，值转为小写 \\
EX（Execution Accuracy） & 预测 SQL 与金标 SQL 在对应 SQLite 库上执行，结果列表完全相等 \\
SQL Validity & 结构校验通过（列索引、聚合、算子、值类型均在合法范围） \\
Execution Success Rate & 结构有效的预测中，可在 SQLite 中成功执行的比例 \\
JSON Validity & 能解析为合法 JSON 对象的比例 \\
Semantic Accuracy（pilot） & 人工判断预测是否与问题意图一致；\texttt{1} 正确、\texttt{0} 错误、\texttt{?} 歧义。本报告仅报告 20 条 pilot，单人标注，未计算 κ \\
\bottomrule
\end{tabular}
\end{minipage}\par\smallskip
EX 使用列表严格相等，而非集合比较；对于多行、无序结果，理论上可能与官方口径产生差异。
\section{实验结果}
\subsection{完整运行结果}
\begin{table*}[t]
\centering
\footnotesize
\caption{四组实验汇总}
\label{tab:t6}
\begin{tabular}{@{}lrrrrrr@{}}
\toprule
实验 & N & LFA (\%) & EX (\%) & JSON Validity (\%) & SQL Validity (\%) & Execution Success (\%) \\
\midrule
E1 & 1,000 & 80.90 & 87.70 & 100.00 & 100.00 & 99.80 \\
E2 & 1,000 & 80.70 & 87.70 & 100.00 & 100.00 & 99.90 \\
E3 & 1,000 & 83.90 & 91.10 & 100.00 & 100.00 & 100.00 \\
\textbf{E4} & \textbf{15,878} & \textbf{85.63} & \textbf{90.77} & \textbf{100.00} & \textbf{100.00} & \textbf{99.94} \\
\bottomrule
\end{tabular}
\end{table*}
E4 的 95\% Wilson 置信区间：LFA [85.07\%, 86.16\%]，EX [90.31\%, 91.21\%]。E4 共 13,596 条 LFA 正确、14,413 条 EX 正确；10 条因算子 \texttt{OP=3} 不可执行。
\begin{quote}\small
注：E1–E3 与 E4 的样本规模不同，表 4 用于汇总运行结果；严格消融结论应以表 5 的同一 1,000 条子集为准。
值得注意的是：E4 在同子集上并不优于 E3，thinking 也无增益；“E4 为主基线”是预设，而非消融选优。
\end{quote}
\subsection{同子集消融}
\begin{table}[t]
\centering
\scriptsize
\caption{相同前 1,000 条 Test 子集上的公平比较}
\label{tab:t7}
\begin{tabular}{@{}lrrrr@{}}
\toprule
实验 & LFA (\%) & EX (\%) & ΔLFA vs E1 (pp) & ΔEX vs E1 (pp) \\
\midrule
E1 & 80.90 & 87.70 & +0.00 & +0.00 \\
E2 & 80.70 & 87.70 & -0.20 & +0.00 \\
E3 & 83.90 & 91.10 & +3.00 & +3.40 \\
E4 & 83.70 & 90.80 & +2.80 & +3.10 \\
\bottomrule
\end{tabular}
\end{table}
\begin{figure*}[t]
\centering
\includegraphics[width=\linewidth]{fig1.pdf}
\caption{0/5-shot 与 thinking 开关的成对消融。误差线为 95\% Wilson 置信区间。}
\end{figure*}
\subsection{成对显著性检验}
\begin{table}[t]
\centering
\footnotesize
\caption{双侧精确 McNemar 检验}
\label{tab:t8}
\begin{tabular}{@{}llrrr@{}}
\toprule
对比 & 指标 & 左对右错 & 左错右对 & p-value \\
\midrule
E1 vs E2 & LFA & 10 & 8 & 0.8145 \\
E1 vs E2 & EX & 8 & 8 & 1.0000 \\
E1 vs E3 & LFA & 13 & 43 & 7.333e-05 \\
E1 vs E3 & EX & 12 & 46 & 8.220e-06 \\
E3 vs E4 & LFA & 10 & 8 & 0.8145 \\
E3 vs E4 & EX & 9 & 6 & 0.6072 \\
\bottomrule
\end{tabular}
\end{table}
结论：5-shot 相对 0-shot 的增益在 LFA 与 EX 上均显著；thinking off/on 在任何一组对比中均不显著。
\subsection{分层结果}
\begin{table}[t]
\centering
\footnotesize
\caption{E4 按聚合类型分层}
\label{tab:t9}
\begin{tabular}{@{}lrrr@{}}
\toprule
聚合 & N & LFA (\%) & EX (\%) \\
\midrule
NONE & 11,324 & 95.10 & 96.35 \\
MAX & 973 & 68.55 & 88.18 \\
MIN & 933 & 70.53 & 89.82 \\
COUNT & 1,455 & 59.18 & 68.25 \\
SUM & 587 & 44.12 & 59.63 \\
AVG & 606 & 63.04 & 76.40 \\
\bottomrule
\end{tabular}
\end{table}
按聚合类型宏平均，LFA 为 \textbf{66.75\%}，EX 为 \textbf{79.77\%}，显著低于微平均的 85.63\% / 90.77\%。这说明整体成绩主要由无聚合查询贡献。
\begin{table}[t]
\centering
\footnotesize
\caption{E4 按 WHERE 条件数量分层}
\label{tab:t10}
\begin{tabular}{@{}lrrr@{}}
\toprule
条件数量 & N & LFA (\%) & EX (\%) \\
\midrule
0 & 131 & 91.60 & 93.13 \\
1 & 10,832 & 87.68 & 92.79 \\
2+ & 4,915 & 80.94 & 86.27 \\
\bottomrule
\end{tabular}
\end{table}
\begin{table}[t]
\centering
\footnotesize
\caption{E4 按表规模、列数、算子与值类型分层（本文新增）}
\label{tab:t11}
\begin{tabular}{@{}llrrr@{}}
\toprule
维度 & 分组 & N & LFA (\%) & EX (\%) \\
\midrule
表行数 & 1–10 & 7,251 & 85.57 & 90.46 \\
表行数 & 11–30 & 7,005 & 85.54 & 90.81 \\
表行数 & 31–100 & 1,388 & 86.24 & 91.93 \\
表行数 & \textgreater{}100 & 234 & 86.32 & 92.74 \\
条件算子 & 含 \texttt{\textgreater{}} 或 \texttt{\textless{}} & 2,928 & 72.10 & 78.52 \\
条件算子 & 全部为 \texttt{=} & 12,950 & 88.69 & 93.54 \\
条件值 & 含数值条件 & 4,073 & 78.76 & 83.75 \\
条件值 & 含字符串条件 & 13,536 & 86.52 & 91.64 \\
列数 & 5–6 列 & 10,292 & 86.05 & 91.06 \\
列数 & 7–9 列 & 4,799 & 85.52 & 90.87 \\
列数 & ≥10 列 & 787 & 80.81 & 86.40 \\
\bottomrule
\end{tabular}
\end{table}
表规模并未导致准确率下降，长表组（\textgreater{}100 行）甚至略高；该组仅 234 条，且长表主要集中在少数表上，不足以得出“长表鲁棒”的结论。真正的成本体现在输入 token 的尾部（最大 132,634）。
\subsection{错误分析}
\begin{table*}[t]
\centering
\footnotesize
\caption{E4 逻辑形式与执行结果交叉表}
\label{tab:t12}
\begin{tabular}{@{}lrrl@{}}
\toprule
类型 & 数量 & 比例 (\%) & 解释 \\
\midrule
LF 正确 / EX 正确 & 13,596 & 85.63 & 逻辑形式与执行结果均正确 \\
LF 错误 / EX 正确 & 817 & 5.15 & 非金标逻辑形式，但在当前表上得到相同答案 \\
LF 正确 / EX 错误 & 0 & 0.00 & 逻辑匹配但执行差异；本轮未出现 \\
LF 错误 / EX 错误 & 1,465 & 9.23 & 真实预测错误 \\
\bottomrule
\end{tabular}
\end{table*}
另有 10 条结构合法但不可执行的预测，全部由 \texttt{OP=3} 导致（0.06\%）。错误高度集中在聚合、非等值条件、数值条件与宽表四类。
\begin{table}[t]
\centering
\footnotesize
\caption{检索 Top-1 分数四分位与性能}
\label{tab:t13}
\begin{tabular}{@{}llrrr@{}}
\toprule
四分位 & 分数范围 & N & LFA (\%) & EX (\%) \\
\midrule
Q1 & 0.239–0.671 & 3,969 & 83.04 & 89.42 \\
Q2 & 0.671–0.776 & 3,970 & 83.35 & 89.29 \\
Q3 & 0.776–0.859 & 3,969 & 86.80 & 91.91 \\
Q4 & 0.859–1.000 & 3,970 & 89.32 & 92.47 \\
\bottomrule
\end{tabular}
\end{table}
检索分数与准确率呈正相关（Q1 → Q4：LFA +6.28 pp，EX +3.05 pp），但这是相关性而非因果证据；高相似度样本本身可能也更简单。
\subsection{效率、缓存与成本}
\begin{figure*}[t]
\centering
\includegraphics[width=\linewidth]{fig2.pdf}
\caption{各组输入/输出 token、缓存命中率与延迟分布。}
\end{figure*}
\subsection{人工语义评测（pilot）}
为检查自动指标是否与真实语义一致，本文对 E4 在 Test 前 20 条上的预测进行了单标注人工评测。评测样本、表格、预测/金标 SQL 及执行结果见 \texttt{reports/\allowbreak{}modern\_baseline/\allowbreak{}human\_eval/\allowbreak{}E4\_semantic\_worksheet.md}，标注记录见 \texttt{E4\_semantic\_labels.csv}。判定标准为：\texttt{1} 表示与问题意图一致且答案等价；\texttt{0} 表示语义错误；\texttt{?} 表示问题歧义或表格信息不足。
\begin{table}[t]
\centering
\footnotesize
\caption{E4 人工语义评测结果（20 条 pilot）}
\label{tab:t14}
\begin{tabular}{@{}lr@{}}
\toprule
指标 & 结果 \\
\midrule
标注样本 & 20 \\
标签 1 / 0 / ? & 20 / 0 / 0 \\
Semantic Accuracy & \textbf{100\%（20/20）} \\
Wilson 95\% CI & \textbf{[83.89\%, 100\%]} \\
同 20 条自动 LFA & 90.00\%（18/20） \\
同 20 条自动 EX & 90.00\%（18/20） \\
标注者数量 & 1 \\
Cohen's κ & 无法计算 \\
\bottomrule
\end{tabular}
\end{table}
人工标签与自动指标完全一致的有 18 条；其余 2 条人工判为语义正确、自动 LFA/EX 判错。没有出现“人工判错、自动判对”的样本。
\textbf{金标错误案例}（\texttt{test:5:1-10083598-1}）：问题为 “What was the number of race that Kevin Curtain won?”。表中 Kevin Curtain 在第 1 行是 Pole Position（杆位），Race winner（分站冠军）为 Kenan Sofuoğlu。预测按 \texttt{Race winner = Kevin Curtain} 计数得 0，金标按 \texttt{Pole Position = Kevin Curtain} 计数得 1。问“赢了几场”时 0 才是正确答案，故该条为金标标注错误，而非模型错误。
\textbf{问题歧义案例}（\texttt{test:19:1-1015521-2}）：问题为 “Give me a list of all spanish officer titles that could receive recognition as group captain in english”。预测选择 Rank in Spanish（\texttt{coronel}），金标选择 Rank in English（\texttt{colonel}）；两种解释均成立，按本工作单规范应标 \texttt{?} 或在备注中记录所选解释。当前标注为 1，因此 100\% 略偏乐观。
\textbf{“实例正确、程序层不可泛化”案例}：S07、S09、S10 都问“有多少个不同的 X”，但预测与金标均使用 \texttt{COUNT(列)} 而非 \texttt{COUNT(DISTINCT)}；由于过滤后恰好只有 1 行，结果正确。这类样本不能说明模型掌握了去重计数语义。
因此，20/20 只能作为 \textbf{pilot 信号}：它说明自动 EX 会低估语义正确的预测（至少 1 例金标错误），但不能作为 E4 在总体上的 Semantic Accuracy；样本为 Test 前 20 条、仅覆盖 6 张表、单人标注且无 κ，尚不具备统计代表性。
\section{讨论}
\subsection{5-shot 检索是主要增益来源}
同子集上 E3 相对 E1 的 LFA/EX 提升为 +3.00/+3.40 pp，且 McNemar p 值均小于 0.001；E3 的执行成功率为 100\%，优于 E1 的 99.80\%。示例帮助模型识别聚合形式、选择列和条件边界，这一点与分层结果一致：5-shot 后 COUNT/SUM 等聚合类错误有所缓解（尽管仍是最弱分支）。
\subsection{Thinking High 未表现出独立收益}
E2 与 E1 的 EX 完全相同，LFA 低 0.20 pp；E4 相对 E3 的差异为 -0.20/-0.30 pp，且不显著。请求中确实记录了 \texttt{reasoning.effort=high}，但 usage 中的 reasoning tokens 并不随配置单调变化（E1 为 4,455，E2 为 5,738，E3 为 6,491，E4 子集为 5,762）。该字段不适合作为推理强度的独立证据，端到端正确率也未显示收益。
\subsection{EX 高于 LFA 的含义}
E4 的 EX 比 LFA 高 5.15 pp，主要来自 817 条“LF 错误 / EX 正确”的等价查询。这些查询虽然条件或选择列与金标不同，但在当前表上返回相同结果。该现象说明：
\begin{itemize}
\item 仅报告 EX 会高估结构生成能力；
\item 仅报告 LFA 会低估实际可用性；
\item 两者应同时报告，并单独分析 denotation equivalence。
\end{itemize}
\subsection{聚合、非等值条件、数值条件是主要短板}
\begin{itemize}
\item SUM 的 LFA 仅 44.12\%，COUNT 为 59.18\%，AVG 为 63.04\%；
\item 含 \texttt{\textgreater{}}/\texttt{\textless{}} 的样本 LFA 为 72.10\%，比全等值条件低 16.59 pp；
\item 含数值条件的样本 LFA 为 78.76\%，比字符串条件低 7.76 pp；
\item ≥10 列的宽表 LFA 为 80.81\%，比 5–6 列低 5.24 pp。
\end{itemize}
这些分支的共同点是需要更精确的列语义匹配与取值规范化，而当前方法主要依赖词法检索和模板化提示词，缺少 schema/value linking 与执行反馈。
\subsection{长表的“准确率不降、成本激增”}
按行数分层，\textgreater{}100 行的表准确率并未下降，但输入 token 尾部极大（P99 12,739，最大 132,634）。这说明模型在本次样本上尚能处理长表，但方法对上下文预算的依赖很强：一旦换用上下文更小的模型，或表行数进一步增长，尾部样本会首先失败。
\subsection{人工语义评测揭示的金标噪声与偶然正确}
20 条 pilot 的语义准确率为 100\%，而同一子集的自动 LFA/EX 为 90\%。差异并非来自模型“多答对了两条”，而是来自评价口径：
\begin{itemize}
\item \textbf{金标噪声}：S06 的问句要求统计胜场，金标却按杆位列计数，预测给出的 0 才是正确答案；
\item \textbf{问题歧义}：S20 的 “spanish officer titles ... in english” 可指西班牙语职称或英文对应职称，人工标注选择了前者；
\item \textbf{偶然正确}：S07/S09/S10 的 COUNT 查询依赖“过滤后只有一行”，并非真正的 DISTINCT 计数，换表即可能失效。
\end{itemize}
这三点共同说明：LFA/EX 不是语义真值，单靠自动指标既可能低估语义正确性，也可能高估程序泛化能力。正式论文应同时报告自动指标、人工语义评测与金标噪声分析；在扩样与双人标注完成前，不应把 20/20 写成总体 Semantic Accuracy。
\section{实验问题}
\subsection{数据集问题}
\subsubsection{SQL 结构高度偏斜}
\needspace{2.5cm}
\noindent\begin{minipage}{\linewidth}
\centering
\scriptsize
\begin{tabular}{@{}>{\raggedright\arraybackslash}p{1.82cm}>{\raggedright\arraybackslash}p{5.58cm}@{}}
\toprule
维度 & 分布 \\
\midrule
聚合 & 无聚合 71.32\%，COUNT 9.16\%，MAX 6.13\%，MIN 5.88\%，AVG 3.82\%，SUM 3.70\% \\
条件数 & 单条件 68.22\%，双条件 24.43\%，三条件及以上 6.52\%，无条件 0.83\% \\
条件算子 & \texttt{=} 82.96\%，\texttt{\textgreater{}} 8.82\%，\texttt{\textless{}} 8.27\% \\
条件值类型 & 字符串 75.19\%，数值 24.81\% \\
\bottomrule
\end{tabular}
\end{minipage}\par\smallskip
这种分布会使系统更容易依赖简单等值查询，训练出的模型通用性不强，也难以用该基准充分验证模型的通用性。
\subsubsection{SQL 语法覆盖过窄}
WikiSQL 的逻辑形式只支持：
\begin{itemize}
\item 单表；
\item 单列 \texttt{SELECT}；
\item 6 种聚合（NONE / MAX / MIN / COUNT / SUM / AVG）；
\item 条件之间只有 AND；
\item 每个条件只有 \texttt{=}、\texttt{\textgreater{}}、\texttt{\textless{}}；
\item 每个条件只有一个值。
\end{itemize}
它不能表达：JOIN、子查询、GROUP BY、HAVING、ORDER BY、LIMIT、DISTINCT、UNION、OR、IN、BETWEEN、LIKE、IS NULL、NOT、算术运算、多列聚合、日期函数、top-k、比例/差值计算。
这导致大量自然语言问句只能被“压扁”成受限形式。对 Test 问句做粗略词面统计：
\needspace{3.6cm}
\noindent\begin{minipage}{\linewidth}
\centering
\scriptsize
\begin{tabular}{@{}>{\raggedright\arraybackslash}p{3.40cm}>{\raggedleft\arraybackslash}p{1.11cm}>{\raggedright\arraybackslash}p{2.89cm}@{}}
\toprule
语言现象 & 问句数（粗略） & WikiSQL 可表达性 \\
\midrule
否定词（not / except / without） & 27 & 无 NOT；常被当成单元格值 \\
or & 73 & 无 OR；通常退化为单条件 \\
between & 7 & 无 BETWEEN；常退化为等值 \\
contains / starts with / ends with & 6 & 无 LIKE；只能等值 \\
top / first / last / highest / lowest / most / least & 1,516 & 部分可由 MAX/MIN 表达，top-k 与排序无法表达 \\
before / after / since / until & 326 & 只有字符串等值与 \texttt{\textgreater{}}/\texttt{\textless{}}，无日期语义 \\
difference / ratio / percentage / times & 154 & 无算术；只能选择已有百分比列 \\
\bottomrule
\end{tabular}
\end{minipage}\par\smallskip
因此，在 WikiSQL 上成功不等于能处理真实业务问题；它更像“表格单元格检索 + 聚合模板匹配”。
\subsubsection{金标错误与问题歧义}
在 20 条人工语义评测样本中发现 2 例问题，其中 1 例为金标错误、1 例为问题歧义：
\begin{itemize}
\item \textbf{金标错误}：\texttt{test:5:1-10083598-1} 问 “What was the number of race that Kevin Curtain won?”。表中 Kevin Curtain 是 Pole Position（杆位），不是 Race winner；金标却按杆位计数得 1，而正确答案是 0。这说明至少部分金标不可靠，不能把 LFA/EX 直接当作语义真值。
\item \textbf{问题歧义}：\texttt{test:19:1-1015521-2} 问 “all spanish officer titles ... in english”，既可理解为输出西班牙语职称（\texttt{coronel}），也可理解为输出英文对应职称（\texttt{colonel}）；金标只选一种解释，且没有歧义标记。
\end{itemize}
\subsubsection{标注噪声与值规范化}
18,114 个等值条件中，有 28 个（0.15\%）即使忽略大小写也找不到完全匹配的单元格。典型问题包括：
\begin{itemize}
\item 逗号/空格不一致：\texttt{Indonesia , Denpasar} vs \texttt{indonesia, denpasar}；
\item 复合值：\texttt{Channing Frye , Jason Richardson (8)} vs \texttt{Channing Frye, Jason Richardson (8)}；
\item 值根本不在表中，或需要跨列/跨行推断。
\end{itemize}
这一比例不高，但说明“字符串等值匹配”并不总是可靠。
\subsubsection{语言质量参差}
问题中存在拼写错误、大小写不一致、缺失标点，例如 “What clu was in toronto 1995-96”“conceeded”“Nejhl” vs “NEJHL”。标注者仍需给出唯一 SQL，解释空间较大。
\subsubsection{列名过于通用}
Test 中最常见表头及出现次数如下：
\needspace{5.3cm}
\noindent\begin{minipage}{\linewidth}
\centering
\footnotesize
\begin{tabular}{@{}lr@{}}
\toprule
表头 & 出现次数 \\
\midrule
\texttt{date} & 1,450 \\
\texttt{score} & 881 \\
\texttt{opponent} & 582 \\
\texttt{record} & 576 \\
\texttt{result} & 524 \\
\texttt{player} & 497 \\
\texttt{year} & 487 \\
\texttt{venue} & 441 \\
\texttt{position} & 433 \\
\texttt{team} & 432 \\
\bottomrule
\end{tabular}
\end{minipage}\par\smallskip
这些列名本身歧义很大，且数据没有列描述、单位说明或业务语义。
\subsubsection{每表覆盖薄与表规模长尾}
\begin{itemize}
\item \textbf{每表问题太少}：Test 共 5,069 张被引用表、15,878 条问题，平均每表约 3.1 条，最多 6 条；单表覆盖很薄，难以支撑表内组合泛化。
\item \textbf{表规模长尾}：最大 1,950 行；75 张表超过 100 行（1.43\%），3 张超过 500 行，3 张超过 1,000 行。长表比例低，但决定了上下文与成本上界；E4 的最大输入达到 132,634 tokens。
\end{itemize}
\subsubsection{答案与语言的双重单一化}
\textbf{微平均掩盖聚合短板。} E4 的微平均 LFA 为 85.63\%，但按聚合类型宏平均只有 66.75\%；SUM 仅 44.12\%，COUNT 59.18\%。这说明模型在真正的聚合推理上远弱于总体数字。
\textbf{答案高度单一。} 金标执行结果的行数分布如下：
\needspace{3.3cm}
\noindent\begin{minipage}{\linewidth}
\centering
\footnotesize
\begin{tabular}{@{}lrr@{}}
\toprule
金标执行结果行数 & 数量 & 占比 \\
\midrule
1 行 & 15,425 & 97.15\% \\
2–5 行 & 366 & 2.31\% \\
6–20 行 & 77 & 0.49\% \\
超过 20 行 & 10 & 0.06\% \\
空结果 & 0 & 0\% \\
\bottomrule
\end{tabular}
\end{minipage}\par\smallskip
也就是说，97\% 的查询只要求返回一个值。这让 EX 很容易饱和：只要选对列、过滤到那一行，结果就正确。
\textbf{条件选择性极高。}
\begin{itemize}
\item 13,638 个等值条件只匹配 1 行；
\item 3,107 个匹配 2–5 行；
\item 1,159 个匹配 6–20 行；
\item 182 个匹配超过 20 行；
\item 28 个匹配 0 行；
\item 11,632 条问题（73.26\%）的所有等值条件唯一确定一行。
\end{itemize}
因此，WikiSQL 主要测的是“单行定位”，而不是多行推理、连接、分组或复杂条件组合。
\textbf{语言高度模板化。}
\begin{itemize}
\item “what is the” 开头 4,320 条（27.2\%）；
\item “what was the” 1,146 条（7.2\%）；
\item 前 15 个三词前缀覆盖 43.2\% 的 Test 问句；
\item Test 中有 28 条问题与 Train 完全相同（表不同），Train–Dev 有 21 条、Dev–Test 有 6 条完全相同。
\end{itemize}
模型可以利用这些表面模式、列位置和常见问法，而不真正理解 schema。
\subsection{可复现与审计问题}
\begin{enumerate}
\item \textbf{Prompt 不可复原}：\texttt{prompt\_freeze.json} 只含版本号、哈希与配置，不含 system/user 正文；\texttt{responses.jsonl} 也不保存实际 Prompt。哈希输入不包括 Prompt 文本、检索权重、示例序列化格式与代码版本。
\item \textbf{版本号不可靠}：\texttt{v2-cache-compact} 与 \texttt{\_v2} 使用同一版本号但哈希不同，版本号未随修改递增。
\item \textbf{人工语义评测规模不足}：目前仅有 Test 前 20 条、单人标注，未计算 Cohen's κ，不能代表总体；\texttt{semantic\_samples.jsonl} 的 \texttt{manual\_label} 仍为 null，标签只存在于 CSV。
\item \textbf{Semantic Accuracy 未进入自动汇总}：\texttt{metrics.json} 的 \texttt{semantic\_accuracy} 仍为 null，实验代码没有读取标注文件的 scorer。
\item \textbf{样本选择偏差}：20 条为 Test 前 20 条，非随机/分层，且以简单等值查询为主，语义准确率被高估。
\end{enumerate}
\section{结论}
本文在完整 WikiSQL Test 集上复核了 gpt-5.6-sol 的 5-shot 基线：\textbf{85.63\% LFA、90.77\% EX}，格式合法性 100\%，执行成功率 99.94\%。同子集消融表明动态 5-shot 检索是主要增益来源（+3.00 pp LFA / +3.40 pp EX，p \textless{} 0.001），而 high reasoning effort 未产生可观测收益。
同时，分层结果揭示了明确的短板：SUM/COUNT 等聚合查询、\texttt{\textgreater{}}/\texttt{\textless{}} 非等值条件、数值条件与宽表上的准确率显著低于总体；按聚合宏平均的 LFA 仅 66.75\%。现有实验在 Dev 验证规模、Test 集复用、提示词冻结证据链、采样参数与成本归因方面存在不足，因此本文将其定位为\textbf{可复核的工程基线}，而非“已完成系统性调优”的最终系统。后续工作应优先补足冻结与复现证据链，再引入更强的检索、schema/value linking 与执行反馈。
人工语义评测 pilot 进一步表明：20/20 的语义正确率与自动 EX 90\% 的差异主要来自 1 条金标错误和 1 条问题歧义，而非模型能力差异；该结果只能作为评价口径的警示，不能替代扩样后的正式语义评测。
\appendices
\needspace{6cm}
\section{数据与代码索引}
\needspace{4.9cm}
\noindent\begin{minipage}{\linewidth}
\centering
\scriptsize
\begin{tabular}{@{}>{\raggedright\arraybackslash}p{3.24cm}>{\raggedright\arraybackslash}p{4.16cm}@{}}
\toprule
内容 & 路径 \\
\midrule
实验代码 & \texttt{experiment/\allowbreak{}run\_modern\_baseline.py} \\
基础模块（数据/API/执行器/指标） & \texttt{experiment/\allowbreak{}run\_experiment.py} \\
正式结果 & \texttt{results/\allowbreak{}formal\_test/\allowbreak{}test/\allowbreak{}E1}–\texttt{E4} \\
主指标 & \texttt{results/\allowbreak{}formal\_test/\allowbreak{}test/\allowbreak{}E4/\allowbreak{}metrics.json} \\
逐样本评测 & \texttt{results/\allowbreak{}formal\_test/\allowbreak{}test/\allowbreak{}E4/\allowbreak{}evaluated.jsonl} \\
原始响应 & \texttt{results/\allowbreak{}formal\_test/\allowbreak{}test/\allowbreak{}E4/\allowbreak{}responses.jsonl} \\
错误用例 & \texttt{results/\allowbreak{}formal\_test/\allowbreak{}test/\allowbreak{}E4/\allowbreak{}error\_cases.md} \\
人工评测工作单 & \texttt{reports/\allowbreak{}modern\_baseline/\allowbreak{}human\_eval/\allowbreak{}E4\_semantic\_worksheet.md} \\
人工标注记录 & \texttt{reports/\allowbreak{}modern\_baseline/\allowbreak{}human\_eval/\allowbreak{}E4\_semantic\_labels.csv} \\
提示词冻结 & \texttt{results/\allowbreak{}formal\_test/\allowbreak{}prompt\_freeze.json} \\
本报告 & \texttt{reports/\allowbreak{}modern\_baseline/\allowbreak{}WikiSQL\_\_Baseline\_\_Report.md} \\
\bottomrule
\end{tabular}
\end{minipage}\par\smallskip
\needspace{6cm}
\section{指标与统计口径}
\begin{itemize}
\item \textbf{LFA}：\texttt{sel}、\texttt{agg} 严格相等；\texttt{conds} 按无序集合比较，值统一小写。重复条件会被集合比较吞掉（金标中罕见）。
\item \textbf{EX}：预测 SQL 与金标 SQL 分别在对应 SQLite 库执行，结果列表严格相等。空结果 \texttt{[]} 视为可执行。
\item \textbf{SQL Validity}：结构校验通过；不保证可执行。
\item \textbf{Execution Success Rate}：结构有效且可执行的样本比例。
\item \textbf{Semantic Accuracy}：人工判断预测是否与问题意图一致；\texttt{1} 正确、\texttt{0} 错误、\texttt{?} 歧义（不计入分母）。当前仅有 20 条 pilot。
\item \textbf{置信区间}：Wilson 区间，95\% 水平。
\item \textbf{McNemar}：双侧精确二项检验，仅使用不一致对。
\item \textbf{宏平均}：对聚合类型六类等权平均，用于抵消 NONE 类的样本量优势。
\end{itemize}
\end{document}
